\documentclass[10pt,a4paper]{amsart}
\usepackage[margin=19mm]{geometry}
\usepackage{amsmath,amssymb,mathtools}
\usepackage[scr=rsfs,cal=euler]{mathalfa}
\usepackage{xfrac}
\usepackage[unicode,hidelinks,bookmarks=false]{hyperref}
\newcommand{\ie}{i.e.\ }
\newcommand{\cf}{cf.\ }
\newcommand{\resp}{respectively\ }
\newcommand{\Subsection}{\subsection}
\providecommand{\nobreakdash}{\nobreak\hbox{-}}
\setlength{\emergencystretch}{2em}
\multlinegap0pt
\usepackage{lmodern}
\usepackage{booktabs}
\usepackage{listings}
\hypersetup{
  pdftitle={A Finite E-Group of Nilpotency Class Three},
  pdfauthor={Xinan Dai, Wenhao Deng, Yingdong Shi, Tailin Wu, Yuchen Yang}
}
\newtheorem{theorem}{Theorem}[section]
\newtheorem{lemma}[theorem]{Lemma}
\newtheorem{proposition}[theorem]{Proposition}
\newtheorem{corollary}[theorem]{Corollary}
\newtheorem{definition}[theorem]{Definition}
\newtheorem{remark}[theorem]{Remark}
\newcommand{\F}{\mathbb F}
\newcommand{\End}{\operatorname{End}}
\newcommand{\Aut}{\operatorname{Aut}}
\newcommand{\im}{\operatorname{im}}
\newcommand{\supp}{\operatorname{supp}}
\newcommand{\rank}{\operatorname{rank}}
\newcommand{\PG}{\operatorname{PG}}
\newcommand{\clq}{\operatorname{cl}_q}
\begin{document}
\title[A Finite E-Group of Nilpotency Class Three]{A Finite E-Group of Nilpotency Class Three}
\author{Xinan Dai, Wenhao Deng, Yidong Shi, Tailin Wu, Yuchen Yang}
\date{}
\maketitle
\noindent\emph{Source and licence.} A Finite E-Group of Nilpotency Class Three. Version arXiv:2608.07275v1 (CC BY 4.0). This material is distributed under the Creative Commons Attribution 4.0 International licence, http://creativecommons.org/licenses/by/4.0/, as recorded in the arXiv OAI licence metadata for this exact version and shown on the abstract page https://arxiv.org/abs/2608.07275v1. Full rights notice: licenses/arxiv-2608.07275-CC-BY-4.0.txt.\par\smallskip
\noindent\textit{Editorial scope.} Counted unit: the original \texttt{theorem} environment \texttt{thm:main} at source lines 595--599 together with its complete proof at lines 601--633; the delivered document keeps source lines 254--699 so that every referenced label and every citation resolves inside it. Native editable source is the paper's own concatenated TeX; the AMS wrapper is editorial formatting only. No publisher class, upstream script or unrelated artwork is redistributed.
\begin{align}
 x_1^3&=[x_2,x_3][x_4,x_5][x_6,x_7][x_8,x_9],\notag\\
 x_2^3&=[x_1,x_3][x_4,x_6][x_5,x_8][x_7,x_9],\notag\\
 x_3^3&=[x_1,x_2][x_4,x_7][x_5,x_9][x_6,x_8],\notag\\
 x_4^3&=[x_1,x_5][x_2,x_6][x_3,x_9][x_7,x_8],\notag\\
 x_5^3&=[x_1,x_4][x_2,x_8][x_3,x_7][x_6,x_9],\label{eq:relations}\\
 x_6^3&=[x_1,x_7][x_2,x_9][x_3,x_5][x_4,x_8],\notag\\
 x_7^3&=[x_1,x_8][x_4,x_9][x_3,x_6][x_2,x_5],\notag\\
 x_8^3&=[x_1,x_9][x_3,x_4][x_2,x_7][x_5,x_6],\notag\\
 x_9^3&=[x_1,x_6][x_3,x_8][x_2,x_4][x_5,x_7].\notag
\end{align}
Each unordered pair $\{i,j\}$ with $1\le i<j\le9$ occurs exactly once on
the right-hand sides.  The presentation and the following structural data
come from \cite[Remark~2.1]{AFH08three} and the consistent presentation in
\cite[p.~2]{AFLObrien10}.

\begin{proposition}[Published structure]\label{prop:published}
The group $P$ has order $3^{84}$ and nilpotency class three.  With
$P_i(P)$ denoting the lower exponent-$3$ central series,
\begin{align}
 \exp(P/P')&=3, & |P/P_2(P)|&=3^{45},\notag\\
 P'=Z_2(P)&\cong C_9^{36}\times C_3^3,
 & \Omega_1(P')=\gamma_3(P)=Z(P)&\cong C_3^{39}.\label{eq:published}
\end{align}
Moreover $P$ is an A-group.
\end{proposition}

Since $P$ is a finite $3$-group,
$\Phi(P)=P^3P'$.  The equality $\exp(P/P')=3$ therefore gives
\begin{equation}\label{eq:phi}
                              \Phi(P)=P'.
\end{equation}
In particular,
\[
                       V:=P/\Phi(P)\cong\F_3^9,
\]
with basis $e_i=x_i\Phi(P)$.  The direct-product description in
\eqref{eq:published} also shows that $P'$ is abelian.

Put $\overline P=P/P_2(P)$.  Because $P_2(P)\leq\Phi(P)$,
$\overline P$ still needs nine generators, so
\[
 |\Phi(\overline P)|=3^{45-9}=3^{36}.
\]
The group $\overline P$ has class at most two, its derived subgroup has
exponent three, and its abelianization has exponent three.  Hence
$\Phi(\overline P)=\overline P'$.  The $36$ commutators
$[x_i,x_j]P_2(P)$ with $i<j$ generate $\overline P'$ and have order at
most three; the order calculation above therefore shows that they form an
$\F_3$-basis.  We obtain a canonical identification, relative to the
chosen basis of $V$,
\begin{equation}\label{eq:wedge-id}
 [x_i,x_j]P_2(P)\longleftrightarrow e_i\wedge e_j,
 \qquad 1\le i<j\le9.
\end{equation}
The class-two computation of Abdollahi--Faghihi--Linton--O'Brien also gives
$|Z(\overline P)|=3^{36}$ \cite[Lemma~3.1 and p.~5]{AFLObrien10}.
Since $\overline P'=\Phi(\overline P)$ already has this order and is
central, we shall use
\begin{equation}\label{eq:centerbar}
 Z(\overline P)=\Phi(\overline P)=\overline P'.
\end{equation}

\section{The relation tensor}\label{sec:tensor}

In a class-two group whose commutator subgroup has exponent three,
$(xy)^3=x^3y^3$.  Thus cubing in $\overline P$ descends to an
$\F_3$-linear map
\begin{equation}\label{eq:q}
                         q:V\longrightarrow\Lambda^2V.
\end{equation}
Using \eqref{eq:wedge-id}, the relations \eqref{eq:relations} give
\begin{align*}
q(e_1)&=e_2\wedge e_3+e_4\wedge e_5+e_6\wedge e_7+e_8\wedge e_9,\\
q(e_2)&=e_1\wedge e_3+e_4\wedge e_6+e_5\wedge e_8+e_7\wedge e_9,\\
q(e_3)&=e_1\wedge e_2+e_4\wedge e_7+e_5\wedge e_9+e_6\wedge e_8,\\
q(e_4)&=e_1\wedge e_5+e_2\wedge e_6+e_3\wedge e_9+e_7\wedge e_8,\\
q(e_5)&=e_1\wedge e_4+e_2\wedge e_8+e_3\wedge e_7+e_6\wedge e_9,\\
q(e_6)&=e_1\wedge e_7+e_2\wedge e_9+e_3\wedge e_5+e_4\wedge e_8,\\
q(e_7)&=e_1\wedge e_8+e_4\wedge e_9+e_3\wedge e_6+e_2\wedge e_5,\\
q(e_8)&=e_1\wedge e_9+e_3\wedge e_4+e_2\wedge e_7+e_5\wedge e_6,\\
q(e_9)&=e_1\wedge e_6+e_3\wedge e_8+e_2\wedge e_4+e_5\wedge e_7.
\end{align*}
This is the sole tensor used below.

\begin{lemma}[Naturality]\label{lem:naturality}
Let $\varphi\in\End(P)$ and let $L\in\End_{\F_3}(V)$ be the induced
linear map.  Then
\begin{equation}\label{eq:naturality}
                   q\circ L=(\Lambda^2L)\circ q.
\end{equation}
Consequently
\begin{equation}\label{eq:imageclosed}
                   q(\im L)\subseteq\Lambda^2(\im L).
\end{equation}
\end{lemma}

\begin{proof}
The subgroup $P_2(P)$ is fully invariant, so $\varphi$ induces an
endomorphism of $\overline P$.  Let $v\in V$ and choose a lift $y\in P$.
Modulo $P_2(P)$ the cube of $y$ represents $q(v)$.  Applying $\varphi$
and using $[a,b]^\varphi=[a^\varphi,b^\varphi]$ gives two descriptions
of the same cube: $q(Lv)$ and $(\Lambda^2L)q(v)$.  This proves
\eqref{eq:naturality}.  Its right-hand side belongs to
$\Lambda^2(\im L)$, which gives \eqref{eq:imageclosed}.
\end{proof}

\begin{definition}\label{def:closed}
A subspace $U\leq V$ is \emph{$q$-closed} if
$q(U)\subseteq\Lambda^2U$.
\end{definition}

The point of Definition~\ref{def:closed} is not to classify all
endomorphisms: Lemma~\ref{lem:naturality} only says that every induced
endomorphism image is $q$-closed.  This one-way implication is exactly
what is needed to rule out singular actions on the Frattini quotient.

\section{Closed subspaces and rigidity}\label{sec:rigidity}

For $\omega\in\Lambda^2V$, contraction identifies $\omega$ with an
alternating map $V^*\to V$.  Define
\[
                    \supp(\omega)=\im(V^*\xrightarrow{\ \omega\ }V).
\]
Equivalently, after choosing the basis $e_1,\ldots,e_9$,
$\supp(\omega)$ is the column space of the skew matrix representing
$\omega$.

\begin{lemma}[Support criterion]\label{lem:support}
For $U\leq V$ and $\omega\in\Lambda^2V$,
\[
                \omega\in\Lambda^2U
       \quad\Longleftrightarrow\quad
                \supp(\omega)\leq U.
\]
\end{lemma}

\begin{proof}
Choose a basis of $U$ and extend it to one of $V$.  If
$\omega\in\Lambda^2U$, its skew matrix has nonzero entries only in the
$U\times U$ block, so its image lies in $U$.  Conversely, if its image
lies in $U$, all rows outside the $U$-block vanish; skew-symmetry forces
the corresponding columns to vanish as well.  Thus no exterior coordinate
of $\omega$ involves a basis vector outside $U$, which is precisely
$\omega\in\Lambda^2U$.
\end{proof}

Starting from $0\neq v\in V$, define $U_0(v)=\langle v\rangle$.  Given
$U_r(v)$, choose a basis $B_r$ and set
\begin{equation}\label{eq:closure}
 U_{r+1}(v)=U_r(v)+\sum_{u\in B_r}\supp(q(u)).
\end{equation}
The process stabilizes because $\dim V=9$.

\begin{lemma}[Least closed subspace]\label{lem:closure}
The stable value $\clq(v)$ of \eqref{eq:closure} is independent of the
chosen bases and is the least $q$-closed subspace containing $v$.
\end{lemma}

\begin{proof}
If $u$ is a linear combination of the vectors in $B_r$, then $q(u)$ is
the same linear combination of their images, and the support of a sum of
bivectors is contained in the sum of the individual supports.  Hence
adjoining the supports for a basis is equivalent to adjoining them for
all vectors in $U_r(v)$; this proves basis independence.

At a fixed point, Lemma~\ref{lem:support} gives
$q(U_r(v))\subseteq\Lambda^2U_r(v)$, so the stable subspace is $q$-closed.
If $W$ is any $q$-closed subspace containing $v$, then
Lemma~\ref{lem:support} shows inductively that $U_r(v)\leq W$ for every
$r$.  Therefore the stable value is the least such subspace.
\end{proof}

\subsection{The exhaustive finite certificate}\label{subsec:certificate}

Scalar multiples have the same closure, so only projective directions
must be checked.  Normalize each nonzero vector by making its first
nonzero coordinate equal to $1$.  The number of normalized vectors is
\begin{equation}\label{eq:pgcount}
                       |\PG(8,3)|=\frac{3^9-1}{3-1}=9841.
\end{equation}
For each normalized $v$, form the alternating matrix $Q(v)$ representing
$q(v)$ and iterate \eqref{eq:closure}, using exact Gaussian elimination
over $\F_3$ for all spans and supports.  The exhaustive calculation gives:

\begin{table}[ht]
\centering
\caption{Exact projective verification for the relation tensor.}
\label{tab:verification}
\begin{tabular}{@{}lrrr@{}}
\toprule
 & rank $4$ & rank $6$ & rank $8$ \\
\midrule
Projective points $[v]$ & 2 & 478 & 9361 \\
Nonzero vectors $v$     & 4 & 956 & 18722 \\
\bottomrule
\end{tabular}

\medskip
\begin{tabular}{@{}lr@{}}
\toprule
Closure statistic & count \\
\midrule
$\dim\clq(v)=9$ & 9841 \\
$\dim\clq(v)<9$ & 0 \\
\bottomrule
\end{tabular}
\end{table}

A finer audit trail records the complete dimension growth of the closure.
Writing, for example, $(1,8,9)$ when the successive dimensions are
$1,8,9$, the exact profile is
\begin{equation}\label{eq:growth-profile}
\begin{array}{c|rrrrrr}
\text{profile}&(1,9)&(1,8,9)&(1,7,9)&(1,6,9)&(1,5,9)&(1,4,9)\\
\hline
\text{count}&6138&3223&426&52&1&1.
\end{array}
\end{equation}
In particular every projective point reaches $V$ after at most two support
enlargements.  The two rank-four directions are represented by
\begin{equation}\label{eq:rank4-points}
 (0,1,1,1,1,1,1,1,1),\qquad
 (1,1,1,1,1,1,1,0,1),
\end{equation}
and have growth profiles $(1,4,9)$ and $(1,5,9)$, respectively.

The first line of Table~\ref{tab:verification} doubles to the published
rank distribution for all nonzero vectors in the corresponding
nine-dimensional relation space \cite[p.~5]{AFLObrien10}.  This is an
independent check that the tensor has been transcribed correctly.  The
closure data are a different computation: they record the iterative
support condition of Lemma~\ref{lem:closure}, not merely the rank of
$q(v)$.

\begin{theorem}[Tensor rigidity]\label{thm:rigidity}
The only $q$-closed subspaces of $V$ are $0$ and $V$.
\end{theorem}

\begin{proof}
Let $U$ be nonzero and $q$-closed.  Choose $0\neq v\in U$.
Lemma~\ref{lem:closure} gives $\clq(v)\leq U$.  By
\eqref{eq:pgcount} and the exhaustive calculation in
Table~\ref{tab:verification}, $\clq(v)=V$ for every nonzero projective
direction.  Hence $U=V$.  The zero subspace is trivially $q$-closed.
\end{proof}

\begin{remark}[What the computation certifies]\label{rem:certificate}
The finite step is exhaustive, not probabilistic: Lemma~\ref{lem:closure}
reduces every nonzero $q$-closed subspace to the closure of one projective
point, and \eqref{eq:pgcount} lists all such points.  A direct exact check
of the $36$ commutator pairs gives $36/36$ distinct pairs.  The reference
verifier in Appendix~\ref{app:verifier} reproduces
Table~\ref{tab:verification} and \eqref{eq:growth-profile} using only
integer arithmetic modulo three.  As a negative control, deleting the
first row $q(e_1)$ leaves $\langle e_1\rangle$ closed, so the same
procedure correctly detects a deliberately introduced proper closed
subspace.
\end{remark}

\subsection{The stabilizer and the orbit question}\label{subsec:stabilizer}

There is an obvious possible shortcut to an exhaustive projective
verification: compute a symmetry group of $q$ and check one point in each
orbit.  In the present example the natural symmetry group is too small for
this strategy.

\begin{proposition}[Trivial projective tensor symmetry]\label{prop:stabilizer}
Set
\[
 \Gamma(q)=\{g\in\operatorname{GL}(V):qg=(\Lambda^2g)q\}
\]
and let
\[
 \Gamma_{\mathrm{sim}}(q)=
 \{g\in\operatorname{GL}(V):qg=\lambda(\Lambda^2g)q
       \text{ for some }\lambda\in\F_3^\times\}.
\]
Then
\[
                 \Gamma(q)=\{I\},\qquad
                 \Gamma_{\mathrm{sim}}(q)=\{I,-I\}.
\]
Consequently the induced action of the natural linear similitude group on
$\PG(V)=\PG(8,3)$ is trivial.
\end{proposition}

\begin{proof}
Let $g\in\Gamma(q)$.  The class-two group $\overline P$ has Frattini
quotient $V$, commutator layer $\Lambda^2V$ via
\eqref{eq:wedge-id}, and power map $q$.  These data give a full class-two presentation of $\overline P$.  Indeed,
introduce central symbols $c_{ij}$ of order three for $1\le i<j\le9$,
impose $[x_i,x_j]=c_{ij}$ and the nine cube relations encoded by $q$, and
kill all triple commutators.  Every element of the resulting group has a
normal form
\[
 x_1^{a_1}\cdots x_9^{a_9}
 \prod_{i<j}c_{ij}^{b_{ij}},\qquad
 0\le a_i,b_{ij}<3,
\]
so its order is at most $3^{9+36}$.  It maps onto $\overline P$, which
has exactly that order; hence the presentation is exact.  The assignments
on $V$ and $\Lambda^2V$ given by $g$ and $\Lambda^2g$ therefore preserve
the defining relations exactly when $qg=(\Lambda^2g)q$.  Thus $g$ lifts
to an endomorphism of $\overline P$.  Since $g$ is invertible on the
Frattini quotient, Burnside's basis theorem makes this lift an
automorphism.

Abdollahi--Faghihi--Linton--O'Brien proved
$\Aut(\overline P)=\Aut_c(\overline P)$
\cite[Lemma~3.1, p.~5]{AFLObrien10}.  By \eqref{eq:centerbar}, a central
automorphism of $\overline P$ acts trivially on
$\overline P/\Phi(\overline P)=V$.  Therefore $g=I$, proving
$\Gamma(q)=\{I\}$.

Now suppose $g\in\Gamma_{\mathrm{sim}}(q)$, so
$qg=\lambda(\Lambda^2g)q$ with
$\lambda\in\F_3^\times=\{1,-1\}$.  Put $h=\lambda g$.  Since
$\lambda^2=1$ and $q$ is linear,
\[
 qh=\lambda qg=\lambda^2(\Lambda^2g)q
    =(\Lambda^2h)q.
\]
Thus $h\in\Gamma(q)$, so $h=I$ and $g=\lambda I$.  Conversely $I$
and $-I$ plainly satisfy the corresponding similitude identities.  Both
scalars induce the identity on projective space.
\end{proof}

\begin{remark}\label{rem:orbit-no-go}
Proposition~\ref{prop:stabilizer} is not needed for the E-group theorem.
Its purpose is to explain the finite proof.  It rules out the most natural
orbit compression: the tensor itself has no nontrivial projective linear
symmetry with which to identify different directions.  It does not rule
out a completely different, non-symmetry-based argument for simplicity of
the dual algebra.  For the present certificate, however, the projective
enumeration is already orbit-minimal for the natural tensor similitude
group.
\end{remark}

\section{Proof of the E-group theorem}\label{sec:main}


\begin{theorem}\label{thm:main}
The group $P$ is a finite E-group of nilpotency class three.  In
particular, Problem~11.46(a) of the Kourovka Notebook has an affirmative
answer.
\end{theorem}

\begin{proof}
Fix $\varphi\in\End(P)$ and let $L$ be its induced map on
$V=P/\Phi(P)$.  By Lemma~\ref{lem:naturality}, $\im L$ is $q$-closed.
Theorem~\ref{thm:rigidity} gives
\begin{equation}\label{eq:dichotomy}
                             \im L=V\quad\text{or}\quad\im L=0.
\end{equation}

Suppose first that $\im L=V$.  Then
$\varphi(P)\Phi(P)=P$.  Burnside's basis theorem implies
$\varphi(P)=P$, and a surjective endomorphism of a finite group is an
automorphism.  Proposition~\ref{prop:published} therefore gives
$[a,a^\varphi]=1$ for every $a\in P$.

Suppose now that $\im L=0$.  By \eqref{eq:phi},
\[
                         \varphi(P)\leq\Phi(P)=P'.
\]
The group $P'$ is abelian.  Apply $\varphi$ to each of the nine relations
in \eqref{eq:relations}.  Every commutator on a right-hand side becomes a
commutator of two elements of $P'$ and hence is trivial.  Therefore
\[
                         \varphi(x_i)^3=1
                         \qquad(1\leq i\leq9).
\]
Since each $\varphi(x_i)$ also lies in $P'$, the identity
$\Omega_1(P')=Z(P)$ from Proposition~\ref{prop:published} yields
$\varphi(x_i)\in Z(P)$ for every $i$.  The $x_i$ generate $P$, so
$\varphi(P)\leq Z(P)$.  Again $[a,a^\varphi]=1$ for every $a\in P$.

The two cases in \eqref{eq:dichotomy} cover all endomorphisms.  The order
and exact nilpotency class are part of Proposition~\ref{prop:published}.
\end{proof}


\begin{remark}\label{rem:notzero}
The second case does not say that an endomorphism inducing zero on
$P/\Phi(P)$ is the zero endomorphism.  It first allows an arbitrary image
inside $P'=Z_2(P)$ and then uses the cube relations to force that image
into the proper subgroup $Z(P)$.  Keeping these two steps separate is
essential in class three.
\end{remark}

\section{Algebraic meaning and scope}\label{sec:meaning}

The condition in Theorem~\ref{thm:rigidity} has a useful dual form.  Let
$A=V^*$ and define a bilinear anticommutative product by
\begin{equation}\label{eq:dualproduct}
 \langle \alpha\cdot\beta,v\rangle
   =\langle \alpha\wedge\beta,q(v)\rangle,
 \qquad \alpha,\beta\in V^*,\ v\in V.
\end{equation}
For $U\leq V$, write
$U^\perp=\{\alpha\in V^*: \alpha(U)=0\}$.

\begin{proposition}\label{prop:dual}
A subspace $U\leq V$ is $q$-closed if and only if $U^\perp$ is an ideal
of the anticommutative algebra $A$.  Consequently Theorem~\ref{thm:rigidity}
is equivalent to simplicity of $A$.
\end{proposition}

\begin{proof}
The subspace $U^\perp$ is an ideal precisely when, for every
$\alpha\in U^\perp$, $\beta\in V^*$, and $u\in U$,
\[
 0=\langle\alpha\cdot\beta,u\rangle
   =\langle\alpha\wedge\beta,q(u)\rangle.
\]
For a fixed $u$, these equalities for all $\alpha\in U^\perp$ and
$\beta\in V^*$ are equivalent to $q(u)\in\Lambda^2U$.  Hence the ideal
condition is equivalent to $q(U)\subseteq\Lambda^2U$.  Orthogonal
complementation reverses $0$ and $V$, so the absence of nonzero proper
$q$-closed subspaces is exactly the absence of nonzero proper ideals in
$A$.
\end{proof}

This reformulation helps locate the rigidity in a broader pattern.  In
Caranti's module-theoretic treatment of class-two groups, power relations
are likewise encoded by a linear map involving $\Lambda^2V$
\cite[Section~3]{Caranti15}; the published erratum should be read together
with that paper for its later modified constructions
\cite{CarantiErratum}.  Caranti's subsequent construction gives examples
where the compatible endomorphisms are reduced to the trivial map and the
identity \cite[Section~4]{Caranti16}.  The present argument has a similar
linear silhouette but a different endpoint: the group is already fixed,
the tensor is the one forced by \eqref{eq:relations}, and a zero action on
the Frattini quotient still has to be lifted through the noncentral group
$P'=Z_2(P)$.

The conclusion therefore rests on two separate rigidity mechanisms.  The
first is linear: the relation tensor admits no nonzero proper closed
subspace, so a singular endomorphism cannot retain any nonzero direction
in $P/\Phi(P)$.  The second is genuinely group-theoretic: once the induced
map vanishes, the nine cube relations collapse the image from $P'$ to
$\Omega_1(P')=Z(P)$.  Neither statement alone proves the E-property.  Their
combination is what makes the class-three candidate work.

\appendix
\section{A self-contained exact reference verifier}\label{app:verifier}


\begin{thebibliography}{99}
\bibitem[AFH08three]{AFH08three} A. Abdollahi, A. Faghihi, and A. Mohammadi Hassanabadi, \emph{3-generator groups whose elements commute with their endomorphic images are abelian}, Comm. Algebra \textbf{36} (2008), no.~10, 3783--3791, \url{https://doi.org/10.1080/00927870802160727}.
\bibitem[AFLObrien10]{AFLObrien10} A. Abdollahi, A. Faghihi, S. A. Linton, and E. A. O'Brien, \emph{Finite $3$-groups of class $3$ whose elements commute with their automorphic images}, Arch. Math. (Basel) \textbf{95} (2010), 1--7, \url{https://doi.org/10.1007/s00013-010-0144-y}.
\bibitem[Caranti15]{Caranti15} A. Caranti, \emph{A module-theoretic approach to abelian automorphism groups}, Israel J. Math. \textbf{205} (2015), no.~1, 235--246, \url{https://doi.org/10.1007/s11856-014-1106-z}.
\bibitem[Caranti16]{Caranti16} A. Caranti, \emph{A simple construction for a class of $p$-groups with all of their automorphisms central}, Rend. Semin. Mat. Univ. Padova \textbf{135} (2016), 251--258, \url{https://doi.org/10.4171/RSMUP/135-14}.
\bibitem[CarantiErratum]{CarantiErratum} A. Caranti, \emph{Erratum to ``A module-theoretic approach to abelian automorphism groups''}, Israel J. Math. \textbf{215} (2016), no.~2, 1025--1026, \url{https://doi.org/10.1007/s11856-016-1401-y}.
\end{thebibliography}
\end{document}
